
  \documentclass[12pt,english]{article}
    \usepackage[T1]{fontenc}
    \usepackage[ansinew]{inputenc}
    %    \usepackage{a4}
    \usepackage{babel}
    \usepackage{fullpage}
 %   \usepackage{german}

\begin{document}
\title{\bf  Undermining the case for Keynesian Demand Management: Is the public expenditure multiplier negative?}
\author{Nikolaus K.A. L�ufer\\
Professor em., University of Konstanz, Germany} \vspace{-1.5mm}
\date{19.8.2009}
\maketitle

\abstract{In one recent note, I have shown that the Haavelmo-Theorem
also holds for open economies. In another recent note, I have
demonstrated that for open economies traditional macroeconomic
analysis is misspecifying either the multiplier or the primary
effect of demand management policies. This is done by neglecting a
parameter for the import content of the components of final demand
shifted. In the present note, the two results are combined to
demonstrate that the classical Haavelmo-Theorem  of a multiplier
equal to 1 is only valid for open economies under a very special
assumption. The assumption is that the import content of public
expenditures is zero. This assumption will hardly ever be met in
modern international and globalised economies. The upshot of this
result is that instead of a multiplier of 1, in open economies we
have in general a multiplier of $1-\frac{\alpha}{1-c+m}$, where
$\alpha$ is the import content parameter for public expenditures
shifted in the experiment. For an import content parameter of a size
equal to 1-c+m, the multiplier becomes zero. Thus, in order to
obtain a zero multiplier an import content parameter of as large as
1 is not required. A much smaller value may be enough to abolish any
kind of multiplier-effect. If the import content parameter is larger
than (1-c+m) then the multiplier may even be negative. Public
expenditure expansion will then contract the economy. Clearly, this
result further undermines the case for Keynesian demand management.}

\section{Introduction}
Keynesian demand management is an idea that continues to fall apart
like an ancient ruin. This happens not because of a negligence from
the side of conserving archaeologists, quite to the contrary, it
follows from the very attempts to conserve the structure by means of
logic and experience. In a prior note I have shown to what extent
the Haavelmo theorem, the first blow to Keynesian economics, can be
extended to open economies. In textbooks, The Haavelmo theorem
usually is demonstrated only within a closed economy framework. In a
second note I have shown that the macroeconomic import function used
in open economy models is grossly erroneous on theoretical grounds.
By combining the two results it may be shown that the Keynesian
expenditure multiplier may easily become close to zero or even
negative. Fortunately, the mathematics involved to demonstrate the
result is rather trivial.


\section{The income model}

We start out by formulating the basic macroeconomic equilibrium
condition:

\begin{equation}
Y = C(Y-T) + A + Ex - Im
\end{equation}
The symbols have the conventional meaning: Y = income, C = private
consumption, T = taxes (lump sum), A = public expenditures, Ex =
exports, Im = imports. For reasons given in a prior note, the import
function to be applied in this case should be:

\begin{equation}
Im = Im(Y-T) = const + \kappa C(Y-T) + \alpha A + \epsilon Ex
\end{equation}
where the conventional marginal propensity to import is
\begin{equation}
m = \kappa \frac{\partial C}{\partial (Y-T) }= \kappa c
\end{equation}
with c being the classical marginal propensity to
consume\footnote{The consumption function $C(Y-T)$ may also be
written as $C_a + C(Y-T)$ where $C_a$ is the absolute, income
independent part of consumption, which sometimes has been called the
existence minimum level of consumption.}, and $\alpha$, $\kappa$,
and $\epsilon$ representing the import content parameters of public
expenditures, private consumption, and exports
respectively. These import content parameters are restricted as x in $0 \leq x \leq 1$.%We
%operate in the context of the Haavelmo Theorem as
%soon as we require
%\begin{equation} dA = dT.
%\end{equation}
\section{The Haavelmo-case}
Additional public expenditures have to be financed either by current
taxes, public debt or by additional money. Restrictions of the
European Monetary Union exclude financing by additional central bank
money. Financing by public debt implies financing by future
taxation. These financing restrictions are appropriately expressed
by the equation
\begin{equation} dA = dT.
\end{equation}
where the tax variable T comprises both current and future taxation
in present value terms.

We need to consider the total differential of the first equation:
\begin{equation}
dY = c(dY-dT) + dA - \alpha dA - m (dY-dT).
\end{equation}
%Here, $m = \kappa c$.

It may be combined with the Haavelmo restriction, $dA=dT$, and
solved for $dY$:
\begin{equation}
dY = (1 - \frac{\alpha}{1-c+m})dA
\end{equation}
If the import content parameter $\alpha$ is zero, then we obtain the
Haavelmo-result of a multiplier of 1, thus replicating the result of
a prior note. However, the multiplier for dA becomes zero if the
following condition holds:
\begin{equation}
\alpha = 1-c + m.
\end{equation}
This condition restricts the value of the import content of public
expenditures, $\alpha$, to the sum of the marginal propensity to
save out of disposable income plus the marginal propensity to import
out of disposable income. The import content parameter $\alpha$ has
not to be equal to 1 in order to obtain a zero public expenditure
multiplier. If $\alpha > 1-c+m$ then the multiplier even becomes
negative. Thus, it is quite possible that public expenditure
expansion has a contractionary effect on the economy.

\section{The case without a tax financing constraint}
If the tax financing constraint is neglected, then dT=0 and
\begin{equation}
dY = \frac{1-\alpha}{1-c+m}dA.
\end{equation}

Since $\alpha \leq 1$, the multiplier cannot be negative in the
present case where the tax financing constraint is neglected. But as
soon as the import content parameter $\alpha$ is larger than $c-m$
the multiplier will be smaller than 1. It even converges to zero as
the import content parameter of additional public demand approaches
1.



\section{Conclusion}
By taking into account the import content of shifts in final demand,
the prospects of Keynesian demand management policies have become
gloomier. The role of the import content parameter can make for
unpleasant surprises in cases where the parameter is neglected,
unknown or where it cannot be known with certainty in advance.  A
good example for this kind of experience is the outcome of the
recent experiment with the used car wrecking premium in Germany. The
statistical data tend to support the view that the action was
counterproductive for the German economy. The role of the import
content parameter (imported cars) was either neglected or
underestimated. To the extent it was neglected this was done in
accordance with current analytical practice of even our best
macroeconomic textbooks.


\bibliographystyle{plain}
\begin{thebibliography}{}
Trygve Haavelmo, Multiplier effects of a Balanced Budget,
Econometrica,
Vol. 13, 1945, S. 311ff. \\
Nikolaus K.A. L�ufer, Haavelmo's theorem in open economies: the
insulation effect, manuscript, May 2009 \\
Nikolaus K.A. L�ufer, Neglecting the import content of shifts in
final demand - a major flaw in current macroeconomic analysis of
open economies, manuscript, August 2009

\end{thebibliography}


\end{document}
